\chapter{Classical Extended-Particle Theories}

\section{Introduction}

The concept of a material point has occupied a central role in theoretical
physics since Newtonian mechanics. In both classical mechanics and
relativistic quantum field theory, elementary particles are commonly idealized
as objects possessing no spatial extension. Mathematically this corresponds to
describing matter by delta-function sources,

\begin{equation}
\rho(\mathbf x)=e\,\delta^{(3)}(\mathbf x),
\end{equation}

or, in covariant notation,

\begin{equation}
J^\mu(x)
=
e
\int
u^\mu(\tau)
\delta^{(4)}
(x-z(\tau))
\,d\tau .
\end{equation}

This idealization is remarkably successful phenomenologically.
Nevertheless, it immediately generates profound conceptual and mathematical
difficulties.

Among them are

\begin{enumerate}
\item infinite electromagnetic self-energy,

\item divergent self-force,

\item runaway solutions,

\item preacceleration,

\item ill-defined radiation reaction,

\item ultraviolet divergences after quantization.
\end{enumerate}

Historically these problems were recognized long before quantum mechanics.
Already by the beginning of the twentieth century it became clear that
Maxwell's equations combined with point charges produce infinite field energy,

\begin{equation}
U
=
\frac{1}{8\pi}
\int
E^2\,d^3x
=
\infty.
\end{equation}

This observation motivated numerous attempts to replace the mathematical point
by an object possessing finite spatial extent.

Although most classical models were eventually abandoned as fundamental
descriptions of elementary particles, many of their central ideas have
reappeared in modern theoretical physics. Examples include nonlinear
electrodynamics, topological solitons, string theory, bag models,
noncommutative geometry, and string-localized quantum fields.

Consequently these early theories are not merely of historical interest.
They represent the first systematic attempts to formulate a relativistic
theory in which elementary particles are extended rather than pointlike.

The purpose of this chapter is to review these classical theories, identify
their achievements and limitations, and explain their lasting influence on
modern relativistic quantum theory.
\section{The Origin of the Point-Particle Problem}

The Coulomb field of a static charge is

\begin{equation}
\mathbf E
=
\frac{e}{r^2}\hat{\mathbf r}.
\end{equation}

The corresponding field energy equals

\begin{equation}
U
=
\frac1{8\pi}
\int
E^2
d^3x
=
\frac{e^2}{2}
\int_0^\infty
\frac{dr}{r^2},
\end{equation}

which diverges at the origin,

\begin{equation}
U=\infty.
\end{equation}

Introducing a finite radius $a$ immediately regularizes the energy,

\begin{equation}
U
=
\frac{e^2}{2a}.
\end{equation}

This simple observation became the principal motivation for extended-electron
models during the period 1900--1930.
\section{Electromagnetic Mass}

Lorentz noticed that the electrostatic energy behaves as inertia.

Using

\begin{equation}
E=mc^2,
\end{equation}

one obtains

\begin{equation}
m_{\rm em}
=
\frac{U}{c^2}
=
\frac{e^2}{2ac^2}.
\end{equation}

This suggested that all of the electron mass might be electromagnetic in
origin.

The idea stimulated an extensive research program aimed at deriving inertia,
momentum and dynamics entirely from Maxwell's equations.\section{Abraham's Rigid Electron}

Max Abraham proposed one of the earliest relativistic extended-particle models.

The electron was assumed to be

\begin{itemize}
\item a rigid sphere,

\item possessing uniform charge density,

\item maintaining constant radius during motion.
\end{itemize}

The electrostatic energy is

\begin{equation}
U
=
\frac35
\frac{e^2}{a},
\end{equation}

for a uniformly charged sphere.

Abraham computed the electromagnetic momentum

\begin{equation}
\mathbf P
=
\frac1{4\pi c}
\int
\mathbf E\times\mathbf B\,d^3x,
\end{equation}

and identified an effective electromagnetic mass.

Although qualitatively successful, the model suffers from a fundamental
difficulty.

Rigid bodies are incompatible with special relativity since perfectly rigid
signals would propagate faster than light.\section{Lorentz's Contracting Electron}

Lorentz proposed a different picture.

Instead of remaining rigid, the electron contracts according to

\begin{equation}
L
=
L_0
\sqrt{1-\frac{v^2}{c^2}}.
\end{equation}

This hypothesis eventually became incorporated into Einstein's special
relativity.

The Lorentz model predicts velocity-dependent longitudinal and transverse
electromagnetic masses,

\begin{align}
m_\parallel
&=
\gamma^3m_0,
\\
m_\perp
&=
\gamma m_0.
\end{align}

Although historically important, these anisotropic masses disappear naturally
within modern relativistic dynamics based upon four-vectors.
\section{Poincaré Stresses}

Electrostatic repulsion tends to make every charged sphere explode.

A stable extended electron therefore requires additional cohesive forces.

Poincaré introduced hypothetical nonelectromagnetic stresses capable of
maintaining equilibrium.

Schematically,

\begin{equation}
T^{\mu\nu}
=
T_{\rm EM}^{\mu\nu}
+
T_{\rm stress}^{\mu\nu}.
\end{equation}

Although successful mathematically, the physical origin of these stresses
remained completely unknown.

This represented perhaps the principal weakness of all classical extended
electron models.\section{The 4/3 Problem}

One of the deepest inconsistencies appeared in the relation between energy and
momentum.

Direct integration gives

\begin{equation}
m_{\rm em}
=
\frac43
\frac{U}{c^2},
\end{equation}

instead of Einstein's expected relation

\begin{equation}
m=\frac{U}{c^2}.
\end{equation}

The discrepancy became known as the
\emph{4/3 problem}.

It was eventually understood that the missing contribution originates from
Poincaré stresses.

In modern relativistic field theory the complete stress-energy tensor satisfies

\begin{equation}
\partial_\mu T^{\mu\nu}=0,
\end{equation}

and the inconsistency disappears once all contributions are included.

\section{Mie's Nonlinear Electrodynamics}

Gustav Mie proposed perhaps the first nonlinear unified field theory.

Instead of Maxwell's linear Lagrangian,

\begin{equation}
L
=
-\frac1{16\pi}
F_{\mu\nu}F^{\mu\nu},
\end{equation}

he considered

\begin{equation}
L
=
L(I_1,I_2),
\end{equation}

where

\begin{align}
I_1
&=
F_{\mu\nu}F^{\mu\nu},
\\
I_2
&=
F_{\mu\nu}\tilde F^{\mu\nu}.
\end{align}

The hope was that nonlinearities would produce finite-energy particle-like
solutions.

Although Mie's particular theory is no longer accepted, the central idea
survives in Born--Infeld electrodynamics and in modern nonlinear gauge
theories.\section{Historical Assessment}

Classical extended-particle theories achieved several lasting successes.

\begin{enumerate}

\item They identified the divergence problem of point charges.

\item They introduced finite-energy field configurations.

\item They emphasized the importance of the stress-energy tensor.

\item They initiated the concept of electromagnetic mass.

\item They inspired nonlinear electrodynamics.

\item They introduced the idea that particles might be localized field
configurations rather than primitive objects.

\end{enumerate}

Their principal shortcomings were equally significant.

\begin{enumerate}

\item The stabilizing forces were ad hoc.

\item No satisfactory relativistic dynamics existed for deformable charged
bodies.

\item Radiation reaction remained unresolved.

\item Spin could not be incorporated naturally.

\item Quantization was unavailable.

Nevertheless, many of their ideas have re-emerged in modern relativistic
theories, often in mathematically refined form.
\begin{thebibliography}{99}

\bibitem{Lorentz1909}
H. A. Lorentz,
\emph{The Theory of Electrons},
Teubner (1909).

\bibitem{Abraham1905}
M. Abraham,
Theorie der Elektrizität,
Vols. I--II,
Teubner (1905).

\bibitem{Poincare1906}
H. Poincaré,
Sur la dynamique de l'électron,
Rendiconti del Circolo Matematico di Palermo
\textbf{21}, 129 (1906).

\bibitem{Mie1912}
G. Mie,
Annalen der Physik
\textbf{342}, 511 (1912).

\bibitem{Rohrlich}
F. Rohrlich,
Classical Charged Particles,
3rd edition,
World Scientific (2007).

\bibitem{Yaghjian}
A. Yaghjian,
Relativistic Dynamics of a Charged Sphere,
Springer (1992).

\end{thebibliography}

\section{Electromagnetic Self-Energy}

Perhaps the oldest indication that the point-particle idealization is
physically problematic is the divergence of the electromagnetic field energy.
Unlike the singularity of the Coulomb potential, which may be regarded merely
as a mathematical inconvenience, the divergence of the total field energy
implies an infinite contribution to the inertial mass of every charged
particle.

This observation predates both quantum mechanics and special relativity.
Already in the last years of the nineteenth century it was realized that the
electromagnetic field itself stores energy and momentum. If the entire mass of
the electron were electromagnetic in origin, then the total field energy
associated with its Coulomb field should equal

\begin{equation}
U = mc^2.
\end{equation}

However, direct calculation immediately produces an ultraviolet divergence.

\subsection{Energy Stored in the Electromagnetic Field}

In Gaussian units the electromagnetic energy density is

\begin{equation}
u
=
\frac{E^2+B^2}{8\pi}.
\end{equation}

For a static charge

\begin{equation}
\mathbf B=0,
\end{equation}

and therefore

\begin{equation}
U
=
\frac1{8\pi}
\int_{\mathbb R^3}
E^2\,d^3x.
\end{equation}

Outside the source the electric field satisfies Maxwell's equations,

\begin{align}
\nabla\cdot\mathbf E &=4\pi\rho,
\\
\nabla\times\mathbf E &=0.
\end{align}

Hence

\begin{equation}
\mathbf E=-\nabla\phi,
\end{equation}

where the electrostatic potential obeys Poisson's equation

\begin{equation}
\nabla^2\phi
=
-4\pi\rho.
\end{equation}

For an ideal point charge

\begin{equation}
\rho(\mathbf r)
=
e\delta^{(3)}(\mathbf r),
\end{equation}

the solution is

\begin{equation}
\phi(r)
=
\frac{e}{r},
\end{equation}

leading to the Coulomb field

\begin{equation}
\mathbf E
=
\frac{e}{r^2}\hat{\mathbf r}.
\end{equation}

Substituting into the energy integral,

\begin{align}
U
&=
\frac1{8\pi}
\int
\frac{e^2}{r^4}
\,d^3x
\\
&=
\frac1{8\pi}
\int_0^\infty
\frac{e^2}{r^4}
4\pi r^2dr
\\
&=
\frac{e^2}{2}
\int_0^\infty
\frac{dr}{r^2}.
\end{align}

The integral diverges at the lower limit,

\begin{equation}
U=\infty.
\end{equation}

Notice that the divergence is entirely ultraviolet. The integral converges
perfectly as $r\rightarrow\infty$,

\begin{equation}
\int_R^\infty
\frac{dr}{r^2}
=
\frac1R,
\end{equation}

whereas

\begin{equation}
\int_0^\varepsilon
\frac{dr}{r^2}
=
\infty.
\end{equation}

Consequently the infinite self-energy originates solely from the assumption
that electric charge occupies zero volume.

\subsection{Regularization by a Finite Radius}

Suppose instead that the charge occupies a sphere of radius $a$.

Outside the sphere the field remains Coulomb,

\begin{equation}
E(r)
=
\frac er^2,
\qquad
r>a.
\end{equation}

If one ignores the internal field for the moment, the exterior energy becomes

\begin{align}
U_{\rm ext}
&=
\frac{e^2}{2}
\int_a^\infty
\frac{dr}{r^2}
\\
&=
\frac{e^2}{2a}.
\end{align}

Immediately the divergence disappears.

The energy is finite for every

\begin{equation}
a>0.
\end{equation}

This elementary observation became the principal motivation for all classical
extended-electron theories.

One should notice that no modification whatsoever of Maxwell's equations has
been required. The divergence results entirely from the singular source
distribution.

\subsection{General Charge Distribution}

The previous calculation may be generalized.

Let

\begin{equation}
\rho(\mathbf r)
\end{equation}

be an arbitrary smooth charge density satisfying

\begin{equation}
\int
\rho(\mathbf r)\,
d^3r
=
e.
\end{equation}

The electrostatic energy may be written in two equivalent forms.

The field representation is

\begin{equation}
U
=
\frac1{8\pi}
\int
E^2\,d^3x.
\end{equation}

Using integration by parts together with Poisson's equation,

\begin{equation}
\nabla^2\phi=-4\pi\rho,
\end{equation}

one obtains

\begin{align}
U
&=
-\frac1{8\pi}
\int
\phi
\nabla^2\phi
\,d^3x
\\
&=
\frac12
\int
\rho\phi
\,d^3x.
\end{align}

Since

\begin{equation}
\phi(\mathbf r)
=
\int
\frac{\rho(\mathbf r')}
{|\mathbf r-\mathbf r'|}
d^3r',
\end{equation}

the energy becomes

\begin{equation}
U
=
\frac12
\int
\int
\frac{\rho(\mathbf r)
\rho(\mathbf r')}
{|\mathbf r-\mathbf r'|}
d^3r\,d^3r'.
\end{equation}

This expression demonstrates immediately that every continuous charge
distribution possesses finite self-energy provided the density remains bounded.

The divergence therefore is not a property of Maxwell's equations but of the
idealization

\begin{equation}
\rho=e\delta^{(3)}(\mathbf r).
\end{equation}

